% ============================================================================
% Fichier  : partie4/P04_C13_incidents_problemes_demandes.tex
% Titre    : Centre de Services, Incidents, Problèmes et Demandes
% Auteur   : MINKA MI NGUIDJOÏ Thierry Emmanuel
% Version  : 1.0 -- octobre 2026
% ============================================================================
% Sources : ITIL 4 (sources publiques, définitions reformulées) ; COBIT 2019,
% fiches DSS02 et DSS03 (dossier Normes). Laboratoire associé : GCI-LAB-06
% (priorisation de tickets, MTTR, taux de réouverture, chez KRIBILOG).
% ============================================================================

\chapter{Centre de Services, Incidents, Problèmes et Demandes}
\label{chap:itil-incidents}
\chaptermark{Incidents, problèmes et demandes}

\epigraph{%
    « Réparer vite, c'est bien.\\
    Ne plus avoir à réparer, c'est mieux. »%
}{Adage des équipes de support}

\niveauChapitre{Intermédiaire}{%
    Maîtrise du Chapitre~\ref{chap:itil-svs} (concepts et système de valeur
    d'ITIL~4).%
}

\begin{boxobjectifs}
\begin{itemize}
    \item Organiser un centre de services.
    \item Prioriser un incident avec une matrice impact $\times$ urgence.
    \item Distinguer incident, problème, erreur connue et demande de
          service.
    \item Calculer les indicateurs (MTTR, taux de résolution au premier
          contact).
    \item Conduire une analyse de cause racine.
\end{itemize}
\end{boxobjectifs}

\bigskip

% ============================================================================
\section*{Question directrice : pourquoi les mêmes pannes reviennent-elles~?}
% ============================================================================

Le Chapitre~\ref{chap:itil-svs} a montré que le flux de valeur
« conteneur bloqué » de KRIBILOG SA s'arrête au contournement~: le
ticket est clos, la cause demeure, et le blocage revient. Avec
4~800~tickets par mois et 30~\% de réouvertures, le centre de services
est saturé, et M.~MBARGA, responsable des services informatiques, ne sait
plus par quoi commencer.

Ce chapitre présente les quatre pratiques qui forment le cœur de
l'exploitation quotidienne~: le centre de services
(section~\ref{sec:c13-desk}), la gestion des incidents
(section~\ref{sec:c13-incidents}), celle des demandes de service
(section~\ref{sec:c13-demandes}) et celle des problèmes
(section~\ref{sec:c13-problemes}). Il apprend à mesurer leur performance
(section~\ref{sec:c13-mesure}) et à les relier à la gouvernance
(section~\ref{sec:c13-gouvernance}). Le laboratoire GCI-LAB-06 fait
prioriser et analyser vingt tickets réels de KRIBILOG.

\separateur

% ============================================================================
\section{Quatre Notions à ne pas Confondre}
\label{sec:c13-notions}
% ============================================================================

\begin{boxdef}
\termecle{Incident}\index{incident}~: interruption non planifiée d'un
service, ou dégradation de sa qualité.

\medskip
\termecle{Problème}\index{problème}~: cause, réelle ou potentielle, d'un ou
de plusieurs incidents.

\medskip
\termecle{Erreur connue}\index{erreur connue} (\textit{known error})~:
problème qui a été analysé, dont la cause est identifiée, mais qui n'est
pas encore résolu.

\medskip
\termecle{Solution de contournement}\index{solution de contournement}
(\textit{workaround})~: moyen de réduire ou d'éliminer l'effet d'un
incident ou d'un problème, sans en traiter la cause.

\medskip
\termecle{Demande de service}\index{demande de service} (\textit{service
request})~: demande d'un utilisateur portant sur une prestation déjà
prévue dans le service (un accès, une information, un équipement).
\end{boxdef}
{\small\itshape Définitions reformulées d'après ITIL~4~\cite{axelos2019itil4}.}

\begin{tableaucours}{L{3.6cm}L{4.6cm}L{4.0cm}}{%
    \tableauHeader{Situation} &
    \tableauHeader{Qualification} &
    \tableauHeader{Pratique}
}
Le logiciel de suivi ne répond plus & Incident & Gestion des incidents \\
Le logiciel perd la position des conteneurs chaque fois que la balise
change de réseau & Problème, puis erreur connue une fois la cause
établie & Gestion des problèmes \\
On relance la synchronisation à la main & Solution de contournement &
Incidents et problèmes \\
Un nouvel agent demande un accès au portail & Demande de service & Gestion
des demandes de service \\
Un agent demande une nouvelle fonction du logiciel & Ni incident ni demande~:
demande de changement & Habilitation des changements
(Chapitre~\ref{chap:itil-changements}) \\
\end{tableaucours}
\vspace{-0.8em}
\begin{center}{\small\itshape Tableau~13.1. Qualifier une situation}\end{center}

La confusion la plus coûteuse est celle de l'incident et du problème. Le
centre de services de KRIBILOG traite chaque blocage comme un incident
isolé~: il rétablit le service, et le blocage revient. Ce qui manque n'est
pas la réactivité, c'est la \emph{gestion des problèmes}.

% ============================================================================
\section{Le Centre de Services}
\label{sec:c13-desk}
% ============================================================================

\subsection{Le point d'entrée unique}

Le \termecle{centre de services}\index{centre de services} (\textit{service
desk}) est le point de contact entre le fournisseur et les utilisateurs
du service. Il reçoit les signalements d'incidents et les demandes, les
enregistre, les qualifie, les résout lorsqu'il le peut, et les oriente
sinon. ITIL~4 insiste sur sa dimension humaine~: il est la vitrine du
service, et c'est par lui que l'utilisateur perçoit la valeur.

\subsection{Organiser le centre de services}

\begin{tableaucours}{L{3.4cm}L{8.8cm}}{%
    \tableauHeader{Choix} &
    \tableauHeader{Options et critères}
}
Canaux & Téléphone, courriel, portail en libre-service, messagerie
instantanée~; au port, le téléphone et la messagerie mobile restent
indispensables \\
Localisation & Centralisé (un seul site), local (près des utilisateurs),
virtuel (équipes réparties, un outil commun) \\
Horaires & Aligner sur l'activité~: le port de Kribi travaille 24~h/24,
le centre de services doit le faire aussi pour les incidents prioritaires \\
Niveaux & Niveau 1 (centre de services), niveau 2 (support spécialisé),
niveau 3 (éditeur, infogérant) \\
Compétences & Écoute, connaissance du métier des utilisateurs, maîtrise
de l'outil et de la base de connaissances \\
\end{tableaucours}
\vspace{-0.8em}
\begin{center}{\small\itshape Tableau~13.2. Les choix d'organisation d'un
centre de services}\end{center}

\begin{boxmetier}{Responsable des services IT}
\textbf{Désengorger le centre de services de KRIBILOG}
\begin{enumerate}
    \item Analyser la répartition des tickets~: 20~\% sont des
          réinitialisations de mot de passe, automatisables par un portail
          en libre-service.
    \item Créer une base de connaissances pour les vingt incidents les plus
          fréquents, avec leur contournement.
    \item Définir des modèles de demandes (accès, équipements) avec leur
          circuit d'approbation.
    \item Exiger, par avenant, que NETPORT SERVICES enregistre toutes ses
          interventions dans l'outil de KRIBILOG.
    \item Nommer un responsable de la gestion des problèmes, distinct du
          responsable du centre de services.
\end{enumerate}
\end{boxmetier}

% ============================================================================
\section{La Gestion des Incidents}
\label{sec:c13-incidents}
% ============================================================================

\subsection{Le cycle d'un incident}

La gestion des incidents a pour but de \emph{rétablir le service le plus
vite possible} et de limiter l'effet de l'incident sur l'activité. Son
cycle comprend l'enregistrement, la classification, la priorisation, le
diagnostic, l'escalade éventuelle, la résolution, puis la clôture après
vérification auprès de l'utilisateur. Ce sont les pratiques DSS02.02 à
DSS02.06 de COBIT (Chapitre~\ref{chap:cobit-objectifs}).

\subsection{La matrice de priorité}

La priorité d'un incident combine deux critères~:
l'\termecle{impact}\index{impact} (l'étendue des effets sur l'activité~:
combien d'utilisateurs, quels processus) et
l'\termecle{urgence}\index{urgence} (la vitesse à laquelle l'effet
s'aggrave). La priorité fixe l'ordre de traitement et le délai cible de
résolution.

\begin{tableaucours}{L{2.6cm}ccc}{%
    \tableauHeader{Impact} &
    \tableauHeader{Urgence élevée} &
    \tableauHeader{Urgence moyenne} &
    \tableauHeader{Urgence faible}
}
Élevé & P1 & P2 & P3 \\
Moyen & P2 & P3 & P4 \\
Faible & P3 & P4 & P5 \\
\end{tableaucours}
\vspace{-0.8em}
\begin{center}{\small\itshape Tableau~13.3. Matrice de priorité impact
$\times$ urgence (exemple)}\end{center}

\begin{tableaucours}{cL{6.0cm}L{4.0cm}}{%
    \tableauHeader{Priorité} &
    \tableauHeader{Exemple à KRIBILOG} &
    \tableauHeader{Délai cible de résolution}
}
P1 & Logiciel de suivi indisponible pour tout le port, navire à quai & 4
heures \\
P2 & Suivi indisponible sur un terminal, chargement en cours & 8 heures \\
P3 & Lenteur générale du logiciel & 24 heures \\
P4 & Imprimante d'un bureau en panne & 72 heures \\
P5 & Erreur d'affichage sans effet sur le travail & 5 jours ouvrés \\
\end{tableaucours}
\vspace{-0.8em}
\begin{center}{\small\itshape Tableau~13.4. Priorités et délais cibles
(convention retenue par KRIBILOG)}\end{center}

Les délais cibles ne sont pas fixés par ITIL~: ils relèvent des
engagements de service convenus avec les métiers
(Chapitre~\ref{chap:itil-niveaux}). Ce que la gouvernance exige, c'est
qu'ils soient \emph{écrits, connus et mesurés}.

\subsection{L'incident majeur}

Un \termecle{incident majeur}\index{incident majeur} est un incident dont
l'impact sur l'activité exige une réponse dépassant le traitement
ordinaire~: une coordination dédiée, une communication spécifique, une
escalade immédiate. La panne de la BANQUE SANAGA
(Chapitre~\ref{chap:prologue}) en est l'exemple type.

\begin{boxprocedure}
\textbf{Le dossier d'incident majeur.} Pour qu'un incident majeur puisse
être analysé, présenté au comité de direction ou au régulateur, le
dossier doit contenir~:
\begin{itemize}
    \item l'heure de détection, l'heure de déclaration en incident majeur
          et le nom de la personne qui l'a déclaré~;
    \item la chronologie des actions, avec leurs auteurs~;
    \item les communications émises (destinataires, heures, contenu)~;
    \item l'heure du rétablissement et la vérification auprès des
          utilisateurs~;
    \item la revue post-incident, avec l'ouverture d'un problème et les
          actions décidées.
\end{itemize}
\end{boxprocedure}

\begin{boxmetier}{DSI}
\textbf{Les trois premières heures d'un incident majeur}
\begin{enumerate}
    \item Désigner immédiatement un responsable de l'incident, seul habilité
          à coordonner, distinct de ceux qui réparent.
    \item Informer la direction générale dans l'heure, avec un point de
          situation régulier, même sans nouvelle.
    \item Décider, à une heure fixée à l'avance, du recours au plan de
          continuité si le service n'est pas rétabli
          (Chapitre~\ref{chap:itil-niveaux}).
    \item Consigner chaque décision~: c'est la chronologie dont le
          régulateur demandera la copie.
\end{enumerate}
\end{boxmetier}

% ============================================================================
\section{La Gestion des Demandes de Service}
\label{sec:c13-demandes}
% ============================================================================

Les demandes de service sont prévisibles~: elles se standardisent. ITIL~4
recommande de définir pour chaque type de demande un modèle (contenu,
circuit d'approbation, délai) et d'offrir autant que possible un
libre-service. La pratique DSS02.03 de COBIT exige de vérifier le droit du
demandeur, d'obtenir les approbations nécessaires et de satisfaire la
demande selon la procédure prévue.

\begin{boxalert}
\textbf{Violation typique.} À l'ANPP, les demandes d'accès au système de
paiement sont traitées par courriel, sans vérification du droit du
demandeur ni approbation de son responsable. Le Chapitre~\ref{chap:audit}
a montré la conséquence~: des comptes ouverts sans justification, et
jamais fermés.

\smallskip
\textbf{Règle~:} toute demande d'accès suit un modèle qui impose
l'approbation du responsable du demandeur et du propriétaire de
l'application, et laisse une trace dans l'outil.
\end{boxalert}

% ============================================================================
\section{La Gestion des Problèmes}
\label{sec:c13-problemes}
% ============================================================================

\subsection{Trois phases}

La gestion des problèmes vise à réduire la probabilité et l'effet des
incidents, en identifiant leurs causes et en gérant les contournements et
les erreurs connues. ITIL~4 la décrit en trois phases~:
\begin{itemize}
    \item l'\textbf{identification des problèmes}~: analyse des tendances
          des incidents, des incidents majeurs, des informations des
          fournisseurs~;
    \item le \textbf{contrôle des problèmes}~: analyse de la cause,
          documentation des contournements et des erreurs connues~;
    \item le \textbf{contrôle des erreurs}~: suivi des erreurs connues
          jusqu'à leur correction, souvent par un changement.
\end{itemize}

\begin{boxobjet}{DSS03 · Gérer les problèmes}
\textbf{Description~:} identifier et classer les problèmes et leurs causes
racines~; les résoudre à temps pour éviter la répétition des incidents~;
formuler des recommandations d'amélioration.

\textbf{Finalité~:} accroître la disponibilité, améliorer les niveaux de
service, réduire les coûts et accroître la satisfaction des clients, en
réduisant le nombre de problèmes d'exploitation.

\textbf{Objectifs liés~:} EG01, EG08~; AG05.

\textbf{Pratiques~:} DSS03.01 identifier et classer les problèmes~;
DSS03.02 analyser et diagnostiquer~; DSS03.03 enregistrer les erreurs
connues~; DSS03.04 résoudre et clore les problèmes~; DSS03.05 conduire une
gestion proactive des problèmes.

\textbf{RACI (extrait, KRIBILOG)~:} DSS03.02~: \emph{A} au responsable des
services informatiques, \emph{R} au responsable des problèmes, \emph{C} à
l'éditeur et à NETPORT SERVICES.

\textbf{Indicateurs~:} nombre de problèmes ouverts et leur âge, part des
incidents liés à une erreur connue, nombre d'incidents évités.
\end{boxobjet}
{\small\itshape Description, finalité, objectifs liés et pratiques résumés
et traduits par nos soins à partir de COBIT~2019~\cite{isaca2018cobitobjectifs}.}

\subsection{L'analyse de cause racine}

\begin{tableaucours}{L{3.2cm}L{4.8cm}L{4.2cm}}{%
    \tableauHeader{Technique} &
    \tableauHeader{Principe} &
    \tableauHeader{Usage}
}
Les cinq pourquoi & Demander « pourquoi~? » jusqu'à la cause (prologue,
section~\ref{sec:pro-causes}) & Causes simples, enchaînées \\
Diagramme d'Ishikawa & Classer les causes possibles par familles
(méthodes, matériel, personnes, environnement\ldots) & Causes multiples \\
Analyse de Pareto & Classer les catégories d'incidents par fréquence & Choisir
les problèmes à traiter en premier \\
Analyse chronologique & Reconstituer la suite des événements & Incidents
majeurs \\
\end{tableaucours}
\vspace{-0.8em}
\begin{center}{\small\itshape Tableau~13.5. Techniques d'analyse de cause}\end{center}

\begin{boxscenario}{Opération KRIBI. Le Pareto des tickets}
\textbf{Faits.} Sur les 4~800~tickets du mois, la répartition est la
suivante~: suivi des conteneurs 1~920 (40~\%), mots de passe 960 (20~\%),
réseau 720 (15~\%), imprimantes 480 (10~\%), divers 720 (15~\%).

\medskip
\textbf{Analyse de Pareto.} Deux catégories, suivi des conteneurs et mots
de passe, représentent 60~\% des tickets. La seconde se traite par le
libre-service. La première justifie l'ouverture d'un problème.

\medskip
\textbf{Cinq pourquoi sur le suivi des conteneurs.} Pourquoi le conteneur
disparaît-il~? Parce que sa position n'est plus mise à jour. Pourquoi~?
Parce que la balise perd la connexion en changeant de réseau mobile entre
le port et l'entrepôt. Pourquoi le logiciel ne la retrouve-t-il pas~?
Parce que la reconnexion n'est pas prévue dans la version installée.
Pourquoi n'a-t-elle pas été corrigée~? Parce que l'éditeur a livré le
correctif il y a six mois, et qu'aucune demande de changement n'a été
ouverte pour l'installer.

\medskip
\textbf{Conclusion.} Erreur connue, contournement documenté (relance
manuelle), correction disponible~: le problème se résout par un
changement (Chapitre~\ref{chap:itil-changements}). Cause de gouvernance~:
personne n'était chargé de suivre les correctifs de l'éditeur.
\end{boxscenario}

% ============================================================================
\section{Mesurer la Performance}
\label{sec:c13-mesure}
% ============================================================================

\subsection{Les indicateurs}

\begin{tableaucours}{L{3.6cm}L{5.0cm}L{3.6cm}}{%
    \tableauHeader{Indicateur} &
    \tableauHeader{Calcul} &
    \tableauHeader{Ce qu'il révèle}
}
Délai moyen de résolution (\sigle{MTTR}) & Somme des durées de résolution
/ nombre d'incidents résolus & Rapidité du rétablissement \\
Résolution au premier contact & Incidents résolus par le niveau~1 sans
escalade / incidents reçus & Efficacité du centre de services \\
Taux de réouverture & Incidents rouverts / incidents clos & Qualité de la
résolution \\
Respect des délais cibles & Incidents résolus dans le délai de leur
priorité / incidents résolus & Tenue des engagements \\
Âge des problèmes ouverts & Durée depuis l'ouverture & Capacité à traiter
les causes \\
\end{tableaucours}
\vspace{-0.8em}
\begin{center}{\small\itshape Tableau~13.6. Indicateurs de la gestion des
incidents et des problèmes}\end{center}

\subsection{Un calcul pas à pas}

\begin{boxscenario}{Opération KRIBI. Huit tickets}
\begin{tableaucours}{ccccc}{%
    \tableauHeader{Ticket} &
    \tableauHeader{Priorité} &
    \tableauHeader{Durée (h)} &
    \tableauHeader{Niveau 1~?} &
    \tableauHeader{Rouvert~?}
}
T1 & P3 & 2 & Oui & Non \\
T2 & P2 & 5 & Non & Oui \\
T3 & P3 & 26 & Non & Non \\
T4 & P4 & 1 & Oui & Non \\
T5 & P1 & 49 & Non & Oui \\
T6 & P2 & 3 & Non & Non \\
T7 & P3 & 8 & Non & Oui \\
T8 & P4 & 30 & Non & Non \\
\end{tableaucours}

\medskip
\textbf{Calculs.} MTTR $= (2+5+26+1+49+3+8+30)/8 = 124/8 = 15{,}5$~heures.
Résolution au premier contact~: $2/8 = 25~\%$. Taux de réouverture~:
$3/8 = 37{,}5~\%$. Respect des délais cibles (Tableau~13.4)~: T3 (26~h
pour P3, cible 24~h) et T5 (49~h pour P1, cible 4~h) sont hors délai, soit
$6/8 = 75~\%$ dans les délais.

\medskip
\textbf{Lecture.} Le MTTR moyen masque l'essentiel~: l'incident P1, le plus
grave, a duré douze fois son délai cible. Un indicateur moyen ne remplace
jamais l'examen des incidents prioritaires.
\end{boxscenario}

\begin{boxalert}
\textbf{Violation typique.} Pour améliorer son taux de résolution,
l'infogérant clôt les tickets dès qu'une solution de contournement est
appliquée, sans vérifier auprès de l'utilisateur. Les indicateurs
s'améliorent, le service se dégrade.

\smallskip
\textbf{Règle~:} tout indicateur doit être accompagné de son indicateur
de contrôle (le taux de réouverture pour le taux de résolution), et la
clôture d'un incident exige la confirmation de l'utilisateur
(DSS02.06).
\end{boxalert}

% ============================================================================
\section{Le Lien avec la Gouvernance}
\label{sec:c13-gouvernance}
% ============================================================================

Les incidents et les problèmes relèvent du management, mais ils
intéressent la gouvernance par trois voies.
\begin{itemize}
    \item \textbf{L'information du conseil}~: les incidents majeurs et les
          tendances doivent remonter au comité des risques ou au comité
          stratégique des SI. Le conseil de la BANQUE SANAGA n'avait reçu
          aucun rapport d'incident en 2025.
    \item \textbf{Le risque}~: un problème non résolu est un risque connu
          et accepté de fait~; il doit figurer au registre des risques
          (APO12) et être accepté explicitement s'il n'est pas traité.
    \item \textbf{Les fournisseurs}~: les incidents imputables à un
          prestataire alimentent sa supervision (APO10) et les pénalités
          prévues au contrat (Chapitre~\ref{chap:itil-niveaux}).
\end{itemize}

\begin{boxmetier}{Auditeur SI}
\textbf{Auditer la gestion des incidents et des problèmes}
\begin{enumerate}
    \item Extraire de l'outil tous les tickets de la période, avec la
          lettre de transmission prévue par le GUID~5100
          (Chapitre~\ref{chap:audit}).
    \item Vérifier sur un échantillon la cohérence de la priorité avec la
          matrice, et le respect des délais cibles.
    \item Rapprocher les tickets rouverts des problèmes ouverts~: une
          catégorie très rouverte sans problème ouvert est un constat.
    \item Vérifier que les incidents majeurs ont donné lieu à une revue
          post-incident.
\end{enumerate}
\end{boxmetier}

% ============================================================================
\separateur
% ============================================================================

\begin{boxresume}
\textbf{Ce qu'il faut retenir de ce chapitre~:}
\begin{itemize}
    \item \textbf{Quatre notions}~: l'incident est une interruption, le
          problème sa cause, l'erreur connue une cause identifiée non
          corrigée, la demande de service une prestation prévue.
    \item \textbf{Le centre de services} est le point d'entrée unique et la
          vitrine du service~; il se désengorge par le libre-service, la
          base de connaissances et les modèles de demandes.
    \item \textbf{La priorité} combine impact et urgence~; les délais cibles
          relèvent des engagements de service et doivent être écrits.
    \item \textbf{L'incident majeur} exige un responsable dédié, une
          communication réglée et un dossier complet.
    \item \textbf{La gestion des problèmes} (DSS03) traite les causes~:
          identification, contrôle des problèmes, contrôle des erreurs.
          Sans elle, les mêmes incidents reviennent.
    \item \textbf{Mesurer}~: MTTR, résolution au premier contact, taux de
          réouverture, respect des délais. Une moyenne ne remplace pas
          l'examen des incidents prioritaires.
\end{itemize}
\end{boxresume}

\bigskip

\begin{boxexo}[title={\ding{45}\quad Exercice~13.1. Qualifier et prioriser \niveaubase}]
Qualifiez chacune des situations suivantes (incident, problème, demande de
service, demande de changement), puis donnez la priorité des incidents
selon la matrice du Tableau~13.3.
\begin{enumerate}[(a)]
    \item Le portail de paiement de l'ANPP est inaccessible la veille de la
          paie des fonctionnaires.
    \item Un agent de la BANQUE SANAGA demande l'installation d'un logiciel
          de bureautique.
    \item Les distributeurs d'une agence de Garoua tombent en panne chaque
          soir vers 22~heures.
    \item Le directeur financier de KRIBILOG ne peut pas imprimer un
          rapport non urgent.
\end{enumerate}
\end{boxexo}

\begin{boxexo}[title={\ding{45}\quad Exercice~13.2. Calculer les indicateurs \niveaumoyen}]
Le centre de services de la BANQUE SANAGA a traité dix incidents dans la
journée, avec les durées de résolution suivantes (en heures)~: 1, 3, 2,
12, 4, 1, 6, 30, 2, 9. Les incidents de 1~heure et le premier incident de
2~heures ont été résolus au premier contact. Deux incidents ont été
rouverts.
\begin{enumerate}[(a)]
    \item Calculez le MTTR, le taux de résolution au premier contact et le
          taux de réouverture.
    \item L'incident de 30~heures était un P1 (cible~: 4~heures). Que
          concluez-vous sur la pertinence du MTTR comme indicateur unique~?
    \item Proposez un tableau de bord de quatre indicateurs pour le comité
          de direction.
\end{enumerate}
\end{boxexo}

\begin{boxexo}[title={\ding{45}\quad Exercice~13.3. La revue post-incident de la BANQUE SANAGA \niveauexpert}]
Conduisez la revue post-incident de la panne de mars~2026
(Chapitre~\ref{chap:prologue}).
\begin{enumerate}[(a)]
    \item Reconstituez la chronologie de l'incident majeur, et indiquez les
          pièces qui manquent au dossier (exigence de contrôle de la
          section~\ref{sec:c13-incidents}).
    \item Ouvrez le problème correspondant et conduisez une analyse
          d'Ishikawa des causes.
    \item Distinguez les actions qui relèvent de la gestion des problèmes,
          de la gestion des changements et de la gouvernance.
    \item Rédigez la partie « incidents » du rapport trimestriel au comité
          des risques.
\end{enumerate}
\end{boxexo}

\bigskip

\begin{boxinfo}
\textbf{Transition.} À KRIBILOG comme à la BANQUE SANAGA, la résolution des
problèmes passe par un changement~: l'installation d'un correctif, la
bascule d'un système. Le Chapitre~\ref{chap:itil-changements} présente
l'habilitation des changements, la gestion des mises en production et la
gestion de la configuration.
\end{boxinfo}
